% TEMPLATE-MILC-v3.tex - plantilla maestra UNICA de las guias MILC v3.
%
% La usan las tres familias de guias, cada una con su driver:
%   · Grados 6-11  -> scripts/build-guias-g11.py
%   · Bebras       -> scripts/build-guias-bebras.py
%   · Semillero    -> scripts/build-guias-semillero.py
%
% Antes habia tres copias casi identicas (~400 lineas iguales, 6 distintas):
% cada mejora habia que aplicarla tres veces, y en la practica no se hacia
% (el motor de recursos vivia solo en dos de ellas). Lo que varia entre
% familias esta parametrizado en 6 placeholders que cada driver llena
% (se nombran aqui SIN su sintaxis real: el builder reemplaza por texto plano,
% tambien dentro de los comentarios, y un valor multilinea rompe el preambulo):
%
%   HEADER_IZQ        encabezado izquierdo de pagina
%   HEADER_DER        encabezado derecho de pagina
%   PORTADA_ETIQUETA  rotulo sobre el numero grande (GRADO/MOMENTO/MODULO)
%   PORTADA_SERIE     linea de serie de la portada
%   PORTADA_ID        identificador de la guia en la portada
%   PORTADA_CIRCULO   el circulito de grado (°), o vacio si no aplica
%
% El resto de claves las documenta content/guias/_SCHEMA.md. El script valida
% que no queden placeholders sin reemplazar antes de escribir el archivo final.

\documentclass[11pt,letterpaper]{article}
\usepackage[margin=2.0cm]{geometry}
\usepackage[table]{xcolor}
\usepackage{fontspec}
\usepackage[spanish]{babel}
\usepackage{tikz}
% Librerías TikZ para los diagramas de las guías (bloque `recursos:`):
%  · babel  → babel en español vuelve activos caracteres como «>», y TikZ los usa
%             en sus opciones (>=stealth, ->). Sin esta librería, cualquier
%             diagrama con flechas revienta con «\language@active@arg> has an
%             extra }». Debe ir ANTES de que se dibuje nada.
%  · positioning        → sintaxis «below left=2pt and 3pt».
%  · decorations.pathreplacing → llaves ({brace}) para señalar rangos.
\usetikzlibrary{babel,positioning,decorations.pathreplacing}
\usepackage{graphicx}
% Circuitos: el trabajo de electrónica se hace hoy con micro:bit y sus sensores
% integrados (MakeCode, sin cableado), así que no se necesitan esquemáticos.
% Si algún día entran montajes con protoboard, basta descomentar esta línea:
% los diagramas .tex/.tikz declarados en `recursos:` ya se incrustan solos.
% \usepackage{circuitikz}
\usepackage[most]{tcolorbox}
\usepackage{tabularx}
\usepackage{array}
\usepackage{fancyhdr}
\usepackage{titlesec}
\usepackage[document]{ragged2e}  % bandera a la derecha en todo el cuerpo
\usepackage{enumitem}
\usepackage{microtype}

% Helvetica Neue no trae la flecha U+2192, asi que cada «→» escrito literalmente
% en el YAML se compilaba a NADA, en silencio (462 flechas en 61 guias: rutas del
% tipo «paleta → tipografia → lockup» salian rotas). Mapeamos el caracter a su
% version matematica, que si tiene glifo. Asi el autor puede seguir escribiendo
% «→» comodamente en el YAML.
\usepackage{newunicodechar}
\newunicodechar{→}{\ensuremath{\rightarrow}}
% Mismo problema con los simbolos de marca y vinetas: el «✓» de «una palomita ✓»
% salia como cuadrito vacio en el PDF de todas las guias que lo usaban.
\usepackage{pifont}
\newunicodechar{✓}{\ding{51}}
\newunicodechar{✗}{\ding{55}}
\newunicodechar{✔}{\ding{52}}
\newunicodechar{•}{\textbullet}
\newunicodechar{≥}{\ensuremath{\geq}}
\newunicodechar{≤}{\ensuremath{\leq}}

\setmainfont{Helvetica Neue}
\setsansfont{Helvetica Neue}
\setlength{\parindent}{0pt}
\setlength{\parskip}{6pt}
% Sin particion de palabras: en 6.º-7.º una palabra cortada por guion al final
% de linea («Comuni-cacion») frena la lectura mucho mas que una linea corta.
\hyphenpenalty=10000
\exhyphenpenalty=10000
\pretolerance=2000
\tolerance=3000
\setlength{\emergencystretch}{3em}
\linespread{1.10}
\renewcommand{\arraystretch}{1.25}
\setlist{leftmargin=5mm,itemsep=1mm,topsep=1mm}
\newcolumntype{Y}{>{\RaggedRight\arraybackslash}X}

\definecolor{milcMagenta}{HTML}{E6007E}
\definecolor{milcVino}{HTML}{5A0038}
\definecolor{milcTurquesa}{HTML}{0093A5}
\definecolor{milcVerde}{HTML}{58A923}
\definecolor{milcMostaza}{HTML}{E5B400}
\definecolor{milcOcre}{HTML}{B86600}
\definecolor{milcNegro}{HTML}{241816}
\definecolor{milcGris}{HTML}{F7F7F4}
\definecolor{milcLinea}{HTML}{E8E2DD}
\definecolor{milcGradePrimary}{HTML}{84CC16}
\definecolor{milcGradeDark}{HTML}{4D7C0F}
\definecolor{milcGradeSoft}{HTML}{ECFCCB}
\definecolor{milcGradeAccent}{HTML}{CFE7FF}
\definecolor{milcGradeText}{HTML}{FFFFFF}
\color{milcNegro}

\pagestyle{fancy}
\fancyhf{}
\lhead{\small Semillero de Investigación · Pensamiento computacional}
\rhead{\small Módulo 1 · Escucha · MILC v3}
\cfoot{\small\thepage}
\renewcommand{\headrulewidth}{0pt}

\newtcolorbox{titlebox}[1]{
  enhanced,colback=#1,colframe=#1,arc=7mm,boxrule=0pt,
  left=7mm,right=7mm,top=3mm,bottom=3mm,
  before skip=8pt,after skip=8pt,
  fontupper=\sffamily\bfseries\Large\color{white}
}

\newtcolorbox{softbox}[3]{
  enhanced,colback=#2,colframe=#1,borderline west={4pt}{0pt}{#1},
  arc=2mm,boxrule=.4pt,left=4mm,right=4mm,top=3mm,bottom=3mm,
  before skip=6pt,after skip=8pt,title={#3},coltitle=white,
  colbacktitle=#1,fonttitle=\sffamily\bfseries,fontupper=\RaggedRight,
  attach boxed title to top left={xshift=3mm,yshift=-2mm},
  boxed title style={arc=2mm, boxrule=0pt}
}

\newtcolorbox{drawbox}[2]{
  enhanced,colback=white,colframe=#1,arc=4mm,boxrule=1pt,
  left=5mm,right=5mm,top=4mm,bottom=4mm,before skip=8pt,after skip=8pt,
  title={#2},coltitle=white,colbacktitle=#1,fonttitle=\sffamily\bfseries,
  fontupper=\RaggedRight,
  attach boxed title to top left={xshift=4mm,yshift=-2mm},
  boxed title style={arc=3mm, boxrule=0pt}
}

\newtcolorbox{infoband}[1]{
  enhanced,colback=#1!8,colframe=#1!50,arc=2mm,boxrule=.5pt,
  left=4mm,right=4mm,top=2mm,bottom=2mm,before skip=4pt,after skip=4pt,
  fontupper=\small\RaggedRight
}

% Puente narrativo entre secciones (contrato editorial Conéctate).
% Da continuidad: "Acabas de X. Ahora vas a Y porque Z."
% Visualmente: banda izquierda en color del grado, texto en itálica pequeña.
\newtcolorbox{puentebox}{
  enhanced,colback=milcGradeSoft,colframe=milcGradeDark,
  borderline west={3pt}{0pt}{milcGradeDark},
  arc=0mm,boxrule=0pt,left=5mm,right=4mm,top=2.5mm,bottom=2.5mm,
  before skip=4pt,after skip=8pt,
  fontupper=\itshape\small\color{milcNegro}\RaggedRight
}
% La flecha va en modo matematico a proposito: Helvetica Neue NO tiene el glifo
% U+2192, asi que \textrightarrow se compilaba a NADA («Missing character: There
% is no → in font Helvetica Neue») y el puente salia sin flecha. En modo
% matematico la flecha viene de la fuente matematica y si se dibuja.
\newcommand{\puente}[1]{%
  \begin{puentebox}%
    \textcolor{milcGradeDark}{$\rightarrow$}\hspace{2mm}#1%
  \end{puentebox}%
}

\newcommand{\linead}[1]{\noindent\color{milcLinea}\rule{#1}{0.5pt}\color{milcNegro}}
\newcommand{\respuesta}[1]{\par\vspace{2mm}\foreach \n in {1,...,#1}{\noindent\color{milcLinea}\rule{\linewidth}{0.5pt}\par\vspace{5mm}}\color{milcNegro}}
\newcommand{\checkbox}{\textcolor{milcGradeDark}{\fbox{\phantom{x}}}\hspace{5pt}}

\newcommand{\phasecard}[4]{
\begin{tcolorbox}[enhanced,colback=#3,colframe=#2,arc=3mm,boxrule=.5pt,height=3.05cm,valign=center,halign=center,left=1.5mm,right=1.5mm,top=2mm,bottom=2mm]
{\sffamily\bfseries\fontsize{22}{24}\selectfont\textcolor{#2}{#1}}\par
{\sffamily\bfseries\small #4}
\end{tcolorbox}
}

% ── Piezas del rediseno 2026 (legibilidad 6.º-11.º) ───────────────────
% Banda de estacion: reemplaza el titlebox macizo. Titulo a la izquierda,
% dato practico (tiempo/modalidad) a la derecha, en una sola linea.
\newcommand{\milcestacion}[2]{%
\begin{tcolorbox}[enhanced,colback=#1,colframe=#1,arc=2mm,boxrule=0pt,
  left=5mm,right=5mm,top=2.6mm,bottom=2.6mm,before skip=11pt,after skip=7pt]
{\sffamily\bfseries\fontsize{15}{18}\selectfont\color{white}#2}
\end{tcolorbox}}

% Tarjeta de paso numerada: sustituye las tablas «Pilar / Aplicacion», que
% eran formato de consulta docente, no de estudiante. El numero va en un
% circulo sobre el borde para que la secuencia se lea de un vistazo.
\newcommand{\milcpaso}[3]{%
\begin{tcolorbox}[enhanced,colback=white,colframe=milcLinea,arc=1.5mm,boxrule=.9pt,
  left=14mm,right=4mm,top=3mm,bottom=3mm,before skip=4pt,after skip=4pt,
  fontupper=\RaggedRight,
  overlay={\node[anchor=north west,circle,fill=#2,text=white,inner sep=0pt,
    minimum size=7.4mm,font=\sffamily\bfseries\small]
    at ([xshift=3.6mm,yshift=-3.2mm]frame.north west) {#1};}]
#3
\end{tcolorbox}}

% Casilla marcable de tamano util con lapiz (la anterior era \fbox{\phantom{x}},
% de alto variable segun el texto que la seguia).
\newcommand{\milcmarca}{\framebox[3.8mm]{\rule{0pt}{3.4mm}}}

% Fila marcable de lista suelta (checks de la estacion 1).
\newcommand{\milccheckrow}[1]{%
\par\vspace{1.8mm}\noindent
\makebox[6.4mm][l]{\raisebox{-.55mm}{\textcolor{milcNegro}{\milcmarca}}}%
\parbox[t]{\dimexpr\linewidth-6.4mm\relax}{\RaggedRight #1}\par}

% Celda marcable dentro de la rubrica (tabla de autoevaluacion). Misma
% sangria francesa, pero pensada para vivir dentro de una columna X.
\newcommand{\milcopcion}[1]{%
\makebox[5.2mm][l]{\raisebox{-.55mm}{\textcolor{milcNegro}{\milcmarca}}}%
\parbox[t]{\dimexpr\linewidth-5.2mm\relax}{\RaggedRight #1}}

% ── Recursos visuales de la guía ──────────────────────────────────────
% Declarados en el bloque `recursos:` del YAML y emitidos por el builder.
% \guiaFigura   → imagen rasterizada o PDF (foto, carta celeste, montaje).
% \guiaDiagrama → fuente TikZ/circuitikz (.tex) incrustada como vector:
%                 es la vía para circuitos y esquemáticos editables.
% #1 = ruta relativa al .tex · #2 = pie (puede ir vacío)
\newcommand{\guiaFigura}[2]{%
\begin{center}
\includegraphics[width=0.88\linewidth,height=0.38\textheight,keepaspectratio]{#1}\\[1.5mm]
{\footnotesize\itshape\textcolor{milcNegro!75}{#2}}
\end{center}
}

\newcommand{\guiaDiagrama}[2]{%
\begin{center}
\input{#1}\\[1.5mm]
{\footnotesize\itshape\textcolor{milcNegro!75}{#2}}
\end{center}
}

\begin{document}
\thispagestyle{empty}

% ── Portada ───────────────────────────────────────────────────────────
% Antes: un rectangulo de color a pagina completa. Se veia bien en pantalla,
% pero en la fotocopiadora del colegio se comia el toner y salia gris sucio.
% Ahora la banda ocupa el tercio superior y el resto va en blanco.
\begin{tikzpicture}[remember picture,overlay]
  \path[fill=milcGradePrimary]
    (current page.north west) rectangle ([yshift=-10.6cm]current page.north east);

  \node[anchor=north west,text=milcGradeText] at ([xshift=2.0cm,yshift=-1.55cm]current page.north west)
    {\sffamily\bfseries\fontsize{12}{14}\selectfont MÓDULO};
  \node[anchor=north west,text=milcGradeText,opacity=.85] at ([xshift=2.0cm,yshift=-2.35cm]current page.north west)
    {\sffamily\bfseries\fontsize{8.5}{10}\selectfont SEMILLERO DE INVESTIGACIÓN · CONECTATE};
  \node[anchor=north east,text=milcGradeText,opacity=.85,align=right] at ([xshift=-2.0cm,yshift=-1.55cm]current page.north east)
    {\sffamily\bfseries\fontsize{9}{12}\selectfont I.E. Sor María Juliana\\Cartago, Valle del Cauca};

  \node[anchor=west,text=milcGradeText] (gradonum) at ([xshift=1.85cm,yshift=-7.1cm]current page.north west)
    {\sffamily\bfseries\fontsize{112}{112}\selectfont 1};
  % El circulito de grado (°) solo aplica a las guias de grado («11°»); en
  % Bebras y Semillero el numero grande es un momento/modulo, no un grado.
  % Se posiciona relativo al nodo `gradonum`, no en coordenadas absolutas.
  

  \node[anchor=west,text=milcGradeText,text width=11.4cm,align=left] at ([xshift=1.5cm,yshift=0.5cm]gradonum.east)
    {
     {\sffamily\bfseries\fontsize{9}{11}\selectfont\MakeUppercase{Módulo 1 · Fase MILC: Escucha · 3 h de clase}}\\[3mm]
     {\sffamily\bfseries\fontsize{21}{25}\selectfont\setlength{\baselineskip}{27pt}Problemas que\\[4pt]piden lógica ---\\[4pt]de la caladora\\[4pt]que cuenta hilos\\[4pt]al problema que\\[4pt]se deja descomponer}\\[3.5mm]
     {\sffamily\bfseries\fontsize{15}{18}\selectfont Pensamiento computacional y algoritmia}
    };
\end{tikzpicture}

\vspace*{9.4cm}
\vfill

{\sffamily\bfseries\fontsize{8}{10}\selectfont\textcolor{milcGradeDark}{LO QUE TE LLEVAS HOY}}

\begin{tcolorbox}[enhanced,colback=white,colframe=milcNegro,arc=2mm,boxrule=1.6pt,
  left=5mm,right=5mm,top=4mm,bottom=4mm,before skip=3pt,after skip=12pt,
  fontupper=\RaggedRight\fontsize{11}{15.5}\selectfont]
Tu \textbf{colección de problemas reformulados}: al menos \textbf{6 problemas cotidianos} de tu casa, tu escuela o tu vereda, cada uno marcado como \emph{problema que pide lógica} o \emph{problema que pide otra cosa}, con la razón. Más \textbf{un problema elegido y descompuesto}: escrito en el formato \textbf{DADO --- SE QUIERE --- TAL QUE}, y abierto en un \textbf{árbol de subproblemas} donde cada rama se pueda resolver por separado. Esa descomposición es la semilla del algoritmo que escribirás en el módulo 2.
\end{tcolorbox}

\begin{tcolorbox}[enhanced,colback=milcGris,colframe=milcGris,arc=2mm,boxrule=0pt,
  left=5mm,right=5mm,top=3.5mm,bottom=3.5mm,after skip=12pt]
{\sffamily\bfseries\fontsize{8}{10}\selectfont\textcolor{milcNegro!55}{ESTA GUÍA ES DE}}\\[3mm]
\begin{tabularx}{\linewidth}{X p{3.2cm} p{3.2cm}}
\linead{\linewidth} & \linead{3.0cm} & \linead{3.0cm}\\[-1mm]
{\fontsize{8.5}{10}\selectfont\textcolor{milcNegro!55}{Nombre completo}} &
{\fontsize{8.5}{10}\selectfont\textcolor{milcNegro!55}{Grupo}} &
{\fontsize{8.5}{10}\selectfont\textcolor{milcNegro!55}{Fecha}}\\
\end{tabularx}
\end{tcolorbox}

\vfill

\noindent\textcolor{milcLinea}{\rule{\linewidth}{1pt}}\\[2mm]
{\sffamily\bfseries\fontsize{9.5}{12}\selectfont Autor: Dr. Álvaro Cárdenas Orozco}\\[1mm]
{\sffamily\fontsize{8.5}{11}\selectfont\textcolor{milcNegro!55}{Metodología MILC · apertura ancestral · pensamiento computacional · triángulo Dussel–estoicismo–Floridi}}

\newpage

\section*{Apertura: Problemas que piden lógica --- de la caladora que cuenta hilos al problema que se deja descomponer}

\begin{softbox}{milcGradeDark}{milcGradeSoft}{Saber ancestral}
\textbf{En Cartago, el calado no se empieza bordando: se empieza contando.} La caladora tensa el lino en el tambor, marca un cuadro sacando \emph{dos hilos horizontales y dos verticales}, y desde ahí retira la urdimbre y la trama \textbf{de uno en uno} ---nunca de a dos---: \emph{``si tratas de deshilar dos hilos a la vez, ellos no te dejan, se pegan entre sí''} (Olivia, caladora de 65 años, en Pérez-Bustos, 2019). Lo que queda no es un hueco: es una \textbf{rejilla regular}, y esa rejilla ---no la tela--- es la base de todo el diseño posterior. Solo entonces la puntada teje hilos nuevos dentro de los agujeros y \emph{repara} lo que se destruyó con cuidado. La señora Elsa, de 80 años, les advirtió a unos ingenieros que fueron a aprender en tres días: lo básico del calado era \textbf{un asunto muy complejo} ---y pasaron cuatro horas en el primer deshilado (Pérez-Bustos, 2019). \textbf{Mira lo que hace su cabeza mientras cuenta:} parte un problema grande en unidades mínimas (hilo por hilo), descubre el patrón que se repite, se queda con la rejilla y suelta el resto, y ejecuta los pasos \emph{en un orden que no se puede alterar}. Eso, con otro nombre, es lo que hoy vas a aprender.
\end{softbox}

\begin{softbox}{milcTurquesa}{milcGris}{Saber contemporáneo}
\textbf{Pensamiento computacional} es formular un problema de modo que su solución se pueda ejecutar \emph{paso a paso}, por una persona, por una máquina o por las dos (Wing, 2006). \textbf{No es programar}: es lo que se piensa \emph{antes} de tocar un teclado ---por eso los retos \textbf{Bebras}, creados por Valentina Dagienė en Lituania, se resuelven con lápiz y papel. Tiene \textbf{cuatro pilares}, y la caladora los usa todos: \textbf{descomposición} (partir el problema en pedazos que se resuelvan solos: el hilo de uno en uno); \textbf{patrones} (ver lo que se repite, para no reinventar cada celda); \textbf{abstracción} (quedarse con lo que decide ---la rejilla--- y soltar el detalle que no decide: el color del lino); \textbf{algoritmo} (una secuencia ordenada, finita y sin ambigüedad, donde \emph{cambiar el orden arruina la pieza}). Y hay algo más que la caladora sabe y el principiante no: \textbf{reconocer cuándo un daño ya no tiene remiendo}. Cuando faltan demasiados hilos, no hay arreglo posible (Pérez-Bustos, 2019). Saber qué problema \emph{sí} se deja resolver ---y cuál no--- también es pensamiento computacional.
\end{softbox}

\begin{softbox}{milcMostaza}{milcGris}{Pregunta-puente}
\textit{La caladora no empieza cortando: empieza \textbf{contando hilos} hasta que el problema se vuelve una rejilla que puede trabajar; \textbf{¿qué problema de tu casa, tu escuela o tu vereda podrías «contar hilo por hilo»} hasta volverlo un reto que se resuelve por pasos\ldots y cuál, por más que cuentes, pide otra cosa?}
\end{softbox}

\begin{softbox}{milcVerde}{milcGris}{Saber y saber hacer}
Al terminar la sesión: (1) podrás \textbf{identificar} problemas de tu entorno que se dejan abordar con pensamiento computacional, y distinguirlos de los que piden otra cosa; (2) sabrás \textbf{explicar} los cuatro pilares ---descomposición, patrones, abstracción, algoritmo--- con un ejemplo propio de cada uno; (3) podrás \textbf{aplicar} el formato DADO --- SE QUIERE --- TAL QUE para reformular un problema cotidiano como reto computacional; (4) habrás \textbf{descompuesto} ese problema en subproblemas que se resuelven por separado.
\end{softbox}



\small
\begin{tabularx}{\textwidth}{>{\bfseries}p{3.0cm}Y}
\rowcolor{milcGradePrimary!12}Autor & Dr. Álvaro Cárdenas Orozco\\
DBA & Formula preguntas investigables, produce evidencia con método y comunica hallazgos reconociendo sus límites (investigación estudiantil STEM, transversal 6°--11°)\\
Referentes & Valentina Dagienė (Bebras) · Jeannette Wing (pensamiento computacional)\\
Producto final & Tu \textbf{colección de problemas reformulados}: al menos \textbf{6 problemas cotidianos} de tu casa, tu escuela o tu vereda, cada uno marcado como \emph{problema que pide lógica} o \emph{problema que pide otra cosa}, con la razón. Más \textbf{un problema elegido y descompuesto}: escrito en el formato \textbf{DADO --- SE QUIERE --- TAL QUE}, y abierto en un \textbf{árbol de subproblemas} donde cada rama se pueda resolver por separado. Esa descomposición es la semilla del algoritmo que escribirás en el módulo 2.\\
\end{tabularx}
\normalsize

\puente{Arranca la línea de Pensamiento computacional. Hoy no vas a programar nada ---ni siquiera vas a encender el computador---: vas a aprender lo que la caladora hace antes de la primera puntada, \textbf{mirar un problema hasta que se deja partir}.}

\milcestacion{milcGradeDark}{Ruta MILC: así vamos a aprender}
\begin{tabularx}{\textwidth}{>{\centering\arraybackslash}X>{\centering\arraybackslash}X>{\centering\arraybackslash}X>{\centering\arraybackslash}X}
\phasecard{1}{milcVerde}{milcGris}{Escucha\\[-1mm]\small Observo, escucho y pregunto.} &
\phasecard{2}{milcTurquesa}{milcGris}{Entender\\[-1mm]\small Sistematizo y ordeno.} &
\phasecard{3}{milcMagenta}{milcGris}{Hacer\\[-1mm]\small Construyo y aplico.} &
\phasecard{4}{milcVino}{milcGris}{Cierre\\[-1mm]\small Evalúo y reflexiono.}
\end{tabularx}


\puente{Antes de resolver, hay que \emph{ver}. Los problemas de tu entorno ya están ahí: la mayoría los pasamos por alto porque nos acostumbramos a que sean así.}

\milcestacion{milcVerde}{Estación 1 de 3 · Escucha}

\begin{softbox}{milcVerde}{milcGris}{Escena para escuchar}
\textbf{Actividad 1 · IDENTIFICA --- Cazador de problemas} (40 min · individual + grupo). \textbf{Momento A (10 min) --- El caso de la caladora.} El profe cuenta cómo se hace un calado en Cartago: marcar el cuadro, sacar los hilos de uno en uno, contar para que la rejilla salga simétrica, remendar el error. Pregunta abierta: \emph{¿qué está haciendo su cabeza mientras cuenta?}. \textbf{Momento B (15 min) --- Safari de problemas.} Cada quien recorre con la mente su casa, su escuela y su vereda buscando \textbf{cosas que se repiten mal, se pierden, se demoran o se enredan}: \emph{``nunca sabemos de quién es el turno de lavar''}, \emph{``se pierde media hora repartiendo los grupos''}, \emph{``la ruta al colegio se llena y no sé cuándo salir''}, \emph{``mi mamá pierde el recibo entre veinte papeles''}. Anota \textbf{mínimo 6}. \textbf{Momento C (15 min) --- ¿Pide lógica o pide otra cosa?} En grupo, separen los problemas que tienen \textbf{datos, reglas y una forma de saber si quedó bien} de los que piden plata, fuerza, permiso o una conversación difícil. \textbf{En el cuaderno:} tus 6+ problemas + la marca de cuáles piden lógica y por qué.
\end{softbox}

{\sffamily\bfseries\fontsize{8}{10}\selectfont\textcolor{milcNegro!55}{MARCA CUANDO LO HAYAS HECHO}}
\milccheckrow{Puedo explicar qué hace la cabeza de la caladora mientras cuenta hilos.}
\milccheckrow{Anoté al menos 6 problemas reales de mi entorno.}
\milccheckrow{Separé los que piden lógica de los que piden otra cosa, y sé decir por qué.}

\begin{infoband}{milcVerde}


\textbf{Lo que va al cuaderno (Actividad 1 · IDENTIFICA).} Título: «Actividad 1 --- Cazador de problemas». \textbf{Formato:} lista de 6+ problemas + a quién le pasa cada uno + la marca «pide lógica / pide otra cosa» con una razón de una línea. \textbf{Extensión:} 12 líneas. \textbf{Sabes que terminaste cuando} puedes señalar \textbf{un problema con datos, reglas y un modo de verificar si quedó bien resuelto}.
\end{infoband}

\puente{Ya tienes problemas cazados. Ahora aprende los cuatro hilos con que se trabajan: descomponer, ver patrones, abstraer y ordenar en un algoritmo.}

\milcestacion{milcTurquesa}{Estación 2 de 3 · Entender}

\begin{softbox}{milcTurquesa}{milcGris}{Lo que vas a entender}
\textbf{Actividad 2 · EXPLICA --- Los cuatro hilos} (65 min · grupo + individual). Ya tienes problemas cazados. Ahora aprende las cuatro herramientas con que se trabajan, viéndolas primero donde llevan siglos funcionando: en las manos de una caladora de Cartago.
\end{softbox}

\milcpaso{1}{milcTurquesa}{\textbf{Pilar 1 --- Descomposición: el hilo de uno en uno.} Un problema grande no se resuelve; se \textbf{parte}. La caladora no arranca un puñado de tela: retira \emph{un hilo a la vez}, porque de a dos \emph{``se pegan entre sí''} y se rompen (Pérez-Bustos, 2019). Descomponer es exactamente eso: cortar el problema por donde \textbf{se deja cortar}, hasta que cada pedazo se pueda resolver \emph{sin mirar a los otros}. Prueba de que cortaste bien: puedes darle un pedazo a un compañero, explicárselo en una frase, y él lo resuelve solo. Si para resolver el pedazo A hay que saber cómo quedó el B, todavía están pegados: \textbf{sigue partiendo}. \emph{Casi todo problema que parece imposible es un montón de problemas fáciles que nadie ha separado.}}
\milcpaso{2}{milcTurquesa}{\textbf{Pilar 2 --- Patrones: lo que se repite no se piensa dos veces.} Cuando la rejilla está lista, la caladora \textbf{no inventa una puntada distinta para cada agujero}: repite la misma forma ---el Punto Espíritu, por ejemplo--- celda tras celda. Reconocer un patrón es descubrir que \textbf{varios pedazos del problema son el mismo pedazo con otro nombre}. Si el turno del lavado, el turno del tablero y el turno del aseo se resuelven con la misma regla de rotación, no tienes tres problemas: tienes \textbf{uno, tres veces}. Buscar patrones ahorra trabajo, sí, pero sobre todo \textbf{revela la estructura}: te dice cuál es de verdad la pieza que hay que resolver. \emph{Quien no ve el patrón termina bordando cuatro horas lo que se resolvía en veinte minutos.}}
\milcpaso{3}{milcTurquesa}{\textbf{Pilar 3 --- Abstracción: quedarse con lo que decide.} La caladora trabaja sobre la \textbf{rejilla}, no sobre la tela: el color del lino, la marca del hilo y el tamaño del tambor no cambian el diseño ---y por eso los suelta. Abstraer es \textbf{quitar todo el detalle que no decide} y quedarse con la estructura que sí. En el problema del turno del lavado: los nombres importan, el orden importa, quién cocinó ayer importa; el color de la loza no. Nombra las \textbf{variables} (lo que cambia), las \textbf{reglas} (lo que siempre se cumple) y el \textbf{criterio de éxito} (cómo sabes que quedó bien). \emph{Abstraer da miedo porque se parece a perder información; en realidad es dejar de cargar la que no sirve.}}
\milcpaso{4}{milcTurquesa}{\textbf{Pilar 4 --- Algoritmo: el orden manda.} Un algoritmo es una \textbf{secuencia ordenada, finita y sin ambigüedad} de pasos que lleva del punto de partida al resultado. El calado tiene el suyo, y es inflexible: marcar el cuadro → sacar dos y dos → deshilar de uno en uno → remendar y bordar. \textbf{Cambia el orden y arruinas la pieza:} bordar antes de deshilar no es «otra manera», es un error irreversible. Un buen algoritmo cumple tres pruebas: cada paso dice \emph{qué} hacer sin dejar dudas, los pasos \textbf{terminan} (no dan vueltas para siempre), y \textbf{otra persona} lo puede ejecutar sin preguntarte nada. Hoy no lo escribes todavía ---eso es el módulo 2---; hoy reconoces que tu problema \textbf{admite} uno.}

\begin{drawbox}{milcTurquesa}{Los cuatro pilares, y el quinto criterio que decide si vale la pena}
\textbf{Descomponer:} partir hasta que cada pedazo se resuelva solo (el hilo de uno en uno). \\[1mm]
\textbf{Patrones:} lo que se repite se resuelve una vez y se reutiliza (la misma puntada en cada celda). \\[1mm]
\textbf{Abstracción:} quedarse con lo que decide ---variables, reglas, criterio de éxito--- y soltar el resto (la rejilla, no el color del lino). \\[1mm]
\textbf{Algoritmo:} pasos ordenados, finitos y sin ambigüedad; cambiar el orden arruina la pieza. \\[1mm]
\textbf{Y antes de todo:} ¿el problema tiene datos, reglas y forma de verificar? Si no, no pide lógica: pide otra cosa.
\end{drawbox}

\begin{softbox}{milcMostaza}{milcGris}{Errores comunes y su corrección}
\textbf{Errores al plantear un problema (y cómo evitarlos).} (1) \emph{Saltar a la solución} --- «hagamos una app»; todavía no sabes cuál es el problema, y la app es solo la primera idea que se te ocurrió. (2) \emph{Subproblemas pegados} --- si para resolver A necesitas el resultado de B y viceversa, no partiste: \textbf{sigue cortando}. (3) \emph{Criterio de éxito invisible} --- «que quede mejor» no es verificable; escribe \textbf{cómo sabrás} que quedó bien. (4) \emph{Confundir difícil con grande} --- casi siempre es grande, no difícil; pártelo. (5) \emph{Abstraer de más} --- si quitas un dato que \emph{sí} decide, tu solución resolverá otro problema. (6) \emph{Empeñarse en lo que no tiene remiendo} --- hay problemas que piden plata, permiso o una conversación; reconocerlo a tiempo es \textbf{criterio}, no rendirse.
\end{softbox}

\begin{infoband}{milcTurquesa}
\guiaDiagrama{assets/pensamiento-computacional-1/rejilla-calado.tex}{Los cuatro pilares no son teoría: son lo que hacen las manos de una caladora de Cartago. Descomponer la tela hilo por hilo deja ver el patrón; quedarse con la rejilla es abstraer; y el orden de los pasos no se puede alterar sin arruinar la pieza. Técnica documentada por Pérez-Bustos (2019).}

\textbf{Lo que va al cuaderno (Actividad 2 · EXPLICA).} Título: «Actividad 2 --- Los cuatro hilos». \textbf{Formato:} 4 bloques, uno por pilar. En cada uno: qué es en tus palabras (2 líneas), cómo lo hace la caladora (1 línea) y \textbf{un ejemplo tuyo} sacado de la Actividad 1 (1 línea). \textbf{Extensión:} 1 hoja. \textbf{Sabes que terminaste cuando} podrías explicarle a alguien \textbf{por qué cambiar el orden de un algoritmo no es «otra manera», sino un error}.
\end{infoband}

\puente{Ya sabes los cuatro pilares. Es hora de \textbf{reformular} un problema tuyo en lenguaje preciso y \textbf{descomponerlo} hasta que cada pedazo se pueda resolver solo.}

\milcestacion{milcMagenta}{Estación 3 de 3 · Hacer}

\begin{softbox}{milcMagenta}{milcGris}{Cómo vas a construir tu producto}
\textbf{Actividad 3 · APLICA --- Reformula y descompone} (45 min · individual + parejas). Hora de trabajar tu propio problema. Vas a escribirlo en \textbf{lenguaje preciso} y a abrirlo en un \textbf{árbol de subproblemas}, hasta que cada rama se sostenga sola.
\end{softbox}

\milcpaso{1}{milcMagenta}{\textbf{Paso 1 --- Elige y reformula (12 min).} De tus problemas que \emph{piden lógica}, escoge el que le importe a una persona concreta. Escríbelo en el \textbf{formato del planteamiento}: \emph{``\textbf{DADO} \_\_\_\_ (los datos y las condiciones que ya existen), \textbf{SE QUIERE} \_\_\_\_ (el resultado, en una sola frase), \textbf{TAL QUE} \_\_\_\_ (las reglas que hay que respetar y cómo sabré que quedó bien)''}. Ejemplo: \emph{``Dado que somos 5 en la casa y cada día hay que lavar, se quiere un turno justo, tal que nadie lave dos días seguidos y todos laven igual número de veces al mes''}.}
\milcpaso{2}{milcMagenta}{\textbf{Paso 2 --- Parte hasta que se suelten (15 min).} Dibuja un \textbf{árbol}: arriba tu problema, y debajo los subproblemas en que se parte. Sigue abriendo ramas mientras la respuesta a \emph{``¿este pedazo se puede resolver sin mirar los otros?''} sea \textbf{no}. Para el turno del lavado: (a) llevar el registro de quién lavó, (b) definir la regla de rotación, (c) resolver los días que alguien no está, (d) avisarle a cada quien su turno. \textbf{Mínimo 3 ramas, y al menos una abierta en dos niveles.}}
\milcpaso{3}{milcMagenta}{\textbf{Paso 3 --- Aplica el «test de la caladora» (10 min).} Rama por rama, pregúntate tres cosas: \emph{¿está sola} (se resuelve sin mirar a las otras)\emph{?}; \emph{¿se puede verificar} (sé cómo comprobar que quedó bien)\emph{?}; \emph{¿tiene remiendo} (o pide plata, permiso o una conversación que no me toca a mí)\emph{?}. Marca cada rama con \textbf{LISTA}, \textbf{PARTIR} (\emph{hay que abrirla más}) o \textbf{SIN REMIENDO} (\emph{esta no se arregla por lógica}). \textbf{No borres las SIN REMIENDO:} nombrar lo que no se puede resolver es parte del trabajo, igual que la caladora que reconoce el daño sin remiendo.}
\milcpaso{4}{milcMagenta}{\textbf{Paso 4 --- Busca el patrón y nómbralo (8 min).} Mira tus ramas y las de tu pareja: \emph{¿hay dos que sean el mismo problema con otro nombre?}. Si «avisarle a cada quien su turno de lavado» y «avisar quién expone el viernes» se resuelven igual, escríbelo: \textbf{«mismo patrón: rotación con aviso»}. \emph{Ese hallazgo es oro: un patrón resuelto una vez sirve muchas veces ---y es justo lo que harás explícito en el módulo 2.}}

\begin{drawbox}{milcMagenta}{Antes de cerrar tu planteamiento}
\checkbox Mi problema está escrito en formato DADO --- SE QUIERE --- TAL QUE. \\[1mm]
\checkbox El criterio de éxito es verificable (sé cómo comprobar que quedó bien). \\[1mm]
\checkbox Mi árbol tiene al menos 3 ramas y una abierta en dos niveles. \\[1mm]
\checkbox Cada rama pasó el test: está sola, se verifica, tiene remiendo. \\[1mm]
\checkbox Marqué SIN REMIENDO lo que no se resuelve por lógica, y no lo borré. \\[1mm]
\checkbox Nombré al menos un patrón que se repite entre ramas.
\end{drawbox}

\begin{softbox}{milcMagenta}{milcGris}{Plantilla de trabajo}
\textbf{Plantilla del planteamiento.} \textit{DADO} \_\_\_\_ (datos y condiciones que ya existen). \textit{SE QUIERE} \_\_\_\_ (resultado, una sola frase). \textit{TAL QUE} \_\_\_\_ (reglas + cómo verifico). \\[1mm]
\textit{Árbol:} problema → subproblema 1 · subproblema 2 · subproblema 3 (abre en dos niveles el que lo necesite). \\[1mm]
\textit{Marca por rama:} LISTA (se sostiene sola) · PARTIR (hay que abrirla más) · SIN REMIENDO (no se resuelve por lógica). \\[1mm]
\textit{Patrón hallado:} «\_\_\_\_ y \_\_\_\_ son el mismo problema: \_\_\_\_».
\end{softbox}

\begin{infoband}{milcMagenta}


\textbf{Lo que va al cuaderno (Actividad 3 · APLICA).} Título: «Actividad 3 --- Mi problema, partido». \textbf{Formato:} el planteamiento DADO/SE QUIERE/TAL QUE + el árbol de subproblemas dibujado + la marca LISTA / PARTIR / SIN REMIENDO de cada rama + el patrón hallado. \textbf{Extensión:} 1 hoja. \textbf{Sabes que terminaste cuando} podrías entregarle \textbf{una rama} a un compañero y él la resolvería sin preguntarte nada más.
\end{infoband}

\puente{El producto no es una solución todavía: es un \textbf{problema bien planteado y bien partido}. La caladora que cuenta mal la rejilla borda cuatro horas para nada.}

\milcestacion{milcMostaza}{Producto de la sesión}
\textbf{Colección de problemas reformulados + un problema descompuesto en árbol (semilla del algoritmo)}.
\begin{softbox}{milcMostaza}{milcGris}{Criterios de calidad}
(1) Hay 6+ problemas del entorno, marcados como «pide lógica» o «pide otra cosa» con su razón. (2) El problema elegido está escrito en formato DADO --- SE QUIERE --- TAL QUE, con criterio verificable. (3) El árbol tiene al menos 3 ramas y una abierta en dos niveles. (4) Cada rama está marcada (LISTA / PARTIR / SIN REMIENDO) y las SIN REMIENDO siguen visibles, no borradas. (5) Se nombra al menos un patrón que se repite entre subproblemas.
\end{softbox}

\puente{Antes de cerrar, revisa tu descomposición: ¿cada subproblema se resuelve sin mirar a los otros, o todavía están pegados como los hilos que salen de a dos?}

\milcestacion{milcVino}{Evaluación liberadora · cinco dimensiones}

\begin{softbox}{milcVino}{milcGris}{Sentido de la evaluación}
Evaluar no es solo revisar una nota. Es reconocer qué aprendí, qué cambió en mí, qué emociones manejé, cómo me relaciono con la ciudad y la comunidad, y qué le dejaré a quien viene detrás.
\end{softbox}

\textbf{Mi reflexión liberadora en cinco dimensiones.} Completa cada frase iniciada:

\small
\begin{tabularx}{\textwidth}{>{\bfseries\RaggedRight\arraybackslash}p{3.4cm}Y>{\RaggedRight\arraybackslash}p{4.6cm}}
\rowcolor{milcVino}\color{white}Dimensión & \color{white}\bfseries Mi respuesta (frame starter) & \color{white}\bfseries Algo que descubrí\\
1. Personal & Lo que más cambió en mí fue \linead{6.5cm} & \linead{4.2cm}\\
2. Emocional & La emoción más fuerte fue \linead{2.5cm} y la regulé \linead{3cm} & \linead{4.2cm}\\
3. Ciudadana & En democracia esto importa porque \linead{6cm} & \linead{4.2cm}\\
4. Local (Cartago) & En mi barrio sería útil para \linead{6.5cm} & \linead{4.2cm}\\
5. Intergeneracional & A mi abuela/abuelo le diría \linead{6.5cm} & \linead{4.2cm}\\
\end{tabularx}
\normalsize

\begin{infoband}{milcVino}
\textbf{Para la próxima sustentación.} Vuelve a esta tabla en seis meses. Verás que las respuestas cambian — no porque lo aprendido cambie, sino porque tú cambias. La evaluación liberadora es una conversación contigo a través del tiempo.
\end{infoband}


\puente{Cerramos con 3 voces que explican por qué plantear bien un problema ---y saber cuál no tiene arreglo--- es la parte difícil, no la fácil.}

\milcestacion{milcGradeDark}{Triángulo de pensamiento · cierre filosófico}

\begin{softbox}{milcVino}{milcGris}{Dussel · lente del nosotros}
\textit{Aquel que es capaz de escuchar silenciosamente al Otro como otro es el que puede constituir una comunidad y no una mera sociedad totalizada.}

Los ingenieros que fueron a Cartago creían que en tres días entenderían «lo básico» del calado, y pasaron cuatro horas sin terminar el primer deshilado (Pérez-Bustos, 2019). \textbf{La caladora sabía algo que ellos no}, y solo escuchándola en silencio ---contando hilos con las manos adoloridas--- pudieron aprenderlo. Frente a un problema pasa igual: quien lo vive sabe cosas que tú no. \textbf{Escuchar antes de plantear no es cortesía: es el único modo de plantear bien.}

\textbf{¿Escuché a quien vive el problema, o lo planteé desde lo que yo creo que le pasa?}
\end{softbox}
\linead{\linewidth}\\[1mm]
\linead{\linewidth}

\begin{softbox}{milcMostaza}{milcGris}{Estoicismo · Epicteto · Enquiridión, 1 (c. 125 d.C.; trad. desde E. Carter) · cuidado interior}
\textit{Algunas cosas están en nuestro poder y otras no. Están en nuestro poder la opinión, el impulso, el deseo, el rechazo; en una palabra, todo aquello que es obra nuestra.}

Epicteto empieza su manual con la operación que hoy practicaste: \textbf{separar lo que depende de ti de lo que no}. Descomponer un problema es eso mismo, con lápiz: esta rama la resuelvo, esta pide plata, esta pide una conversación que no me toca. \textbf{Marcar una rama «sin remiendo» no es rendirse ---es dejar de gastar fuerza donde no hay arreglo}, para ponerla toda donde sí. El que no separa termina agotado y sin resolver nada.

\textbf{De todo lo que me molesta de este problema, ¿qué parte está de verdad en mi poder resolver?}
\end{softbox}
\linead{\linewidth}\\[1mm]
\linead{\linewidth}

\begin{softbox}{milcTurquesa}{milcGris}{Floridi · lente de la infoesfera}
\textit{Los pequeños patrones solo pueden ser significativos si se agregan correctamente, se comparan y se procesan a tiempo.}

Floridi advierte que un patrón suelto no dice nada: adquiere sentido cuando se \textbf{junta con otros, se compara y llega a tiempo}. Una puntada aislada no es un calado; la misma puntada repetida en la rejilla \emph{sí} es un diseño. Por eso el paso 4 de hoy ---\textbf{buscar qué ramas son el mismo problema con otro nombre}--- no es un adorno: es lo que convierte tu montón de subproblemas en una estructura que puede resolverse una vez y reutilizarse muchas.

\textbf{¿Qué patrón se repite entre mis subproblemas, y qué me ahorraría resolverlo una sola vez bien?}
\end{softbox}
\linead{\linewidth}\\[1mm]
\linead{\linewidth}

\begin{infoband}{milcNegro}
\textbf{Nota para el docente.} Las tres formulaciones del triángulo son la síntesis del autor de la guía, no citas textuales de los pensadores. Si vas a citarlos en clase, remite a la obra original.
\end{infoband}

\milcestacion{milcVino}{Mi autoevaluación · marca una casilla por fila}

{\small
\setlength{\extrarowheight}{2.2pt}
\begin{tabularx}{\textwidth}{>{\RaggedRight\arraybackslash\bfseries}p{3.0cm}YYY}
\rowcolor{milcVino}\color{white}\bfseries Criterio & \color{white}\bfseries Lo logré & \color{white}\bfseries En proceso & \color{white}\bfseries Necesito apoyo\\[.6mm]
Comprensión del tema & \milcopcion{Explico las ideas con mis palabras.} & \milcopcion{Reconozco las ideas, falta precisión.} & \milcopcion{Necesito repasar lo básico.}\\[.8mm]
\rowcolor{milcGris}Pensamiento computacional & \milcopcion{Usé los pilares al armar mi producto.} & \milcopcion{Usé algunos pilares.} & \milcopcion{Necesito releer la estación 2.}\\[.8mm]
Aplicación y producto & \milcopcion{Mi producto cumple los criterios.} & \milcopcion{Está armado, con ajustes pendientes.} & \milcopcion{Aún está en construcción.}\\[.8mm]
\rowcolor{milcGris}Reflexión 5 dimensiones & \milcopcion{Respondí cada una con honestidad.} & \milcopcion{Respondí algunas con detalle.} & \milcopcion{Mis respuestas son muy generales.}\\[.8mm]
Triángulo de pensamiento & \milcopcion{Conecté las 3 voces con mi trabajo.} & \milcopcion{Conecté al menos una.} & \milcopcion{No las relaciono con mi proceso.}\\[.8mm]
\end{tabularx}}

\vspace{3mm}
\textbf{Hoy entendí que:}\\[1mm]
\linead{\linewidth}\\[1mm]
\linead{\linewidth}

\textbf{Mi compromiso para la próxima sesión es:}\\[1mm]
\linead{\linewidth}\\[1mm]
\linead{\linewidth}

\begin{softbox}{milcMostaza}{milcGris}{Cita de despedida · Marco Aurelio}
\textit{``El obstáculo en el camino se vuelve el camino. Lo que estorba al camino se vuelve camino.''}\\[1mm]
Cada error de hoy, bien observado, se convierte en pieza del camino que pisarás mañana.
\end{softbox}

\small
\begin{tabularx}{\textwidth}{>{\bfseries}p{3.6cm}Y}
\rowcolor{milcGradeDark!12}Nombre completo: & \linead{10cm}\\
Fecha: & \linead{4cm} \quad Curso/grupo: \linead{4cm}\\
Mi firma: & \linead{9cm}\\
\end{tabularx}
\normalsize

\begin{infoband}{milcGradeDark}
\textbf{Para la próxima sesión.} Dos cosas. \textbf{(1) Levanta el algoritmo de un oficio:} busca a alguien de tu casa o tu barrio que haga algo con pasos ---calar o bordar, amasar el pan, arreglar una bicicleta, preparar el sancocho, cuadrar una caja--- y pídele que te lo explique mientras lo hace. Escríbelo en \textbf{6 a 10 pasos numerados} y luego pregúntale: \emph{¿qué pasa si se cambia el orden de estos dos?}. Anota su respuesta \textbf{tal como la diga}: ahí está la diferencia entre un paso que se puede mover y uno que no. \textbf{(2) Valida tu planteamiento con su dueño:} muéstrale a la persona a quien le importa tu problema el formato DADO --- SE QUIERE --- TAL QUE y pregúntale si eso \emph{es} su problema. Si te corrige, ganaste. \textbf{Trae:} el algoritmo del oficio y tu planteamiento corregido. \emph{En el módulo 2 vas a convertir tu árbol en patrones y en un algoritmo escrito ---diagrama de flujo y pseudocódigo---, y en el 3 lo pondrás a prueba contra los retos Bebras. La caladora ya te enseñó lo más difícil: contar antes de cortar.}
\end{infoband}

\end{document}
